\section{Method}
\label{sec:method}

Keep paper-specific notation in \texttt{main.tex}. Displayed equations are numbered and can be referenced with \verb|\cref|, as in \cref{eq:equivariance}:
\begin{equation}
  f(\rho_{\mathrm{in}}(g)\,x) = \rho_{\mathrm{out}}(g)\,f(x), \qquad \forall g \in G.
  \label{eq:equivariance}
\end{equation}

\newtheorem{theorem}{Theorem}
\begin{theorem}\label{thm:main}
Every layer of \method commutes with the action of $\mathrm{E}(3)$, hence so does the composed network.
\end{theorem}

\begin{algorithm}[t]
  \caption{Training loop of \method}\label{alg:train}
  \begin{algorithmic}[1]
    \Require dataset $\mathcal{D}$, learning rate $\eta$
    \For{each minibatch $(x, y) \in \mathcal{D}$}
      \State $\hat{y} \gets f_\theta(x)$
      \State $\theta \gets \theta - \eta\,\nabla_\theta \ell(\hat{y}, y)$
    \EndFor
  \end{algorithmic}
\end{algorithm}

\paragraph{Implementation} Run-in paragraph headings use \verb|\paragraph| or \verb|\parahead|. \Cref{alg:train} shows the algorithm environment.
